% !TeX root = ../../tesis.tex
%%%%%%%%%%%%%%%%%%%%%%%%%%%%%%%%%%%%%%%%%%%%%%%%%%%%%%%%%%%%%%%%%%%
% SECCIÓN 5.3.3 — FUERZA CAPILAR (FC)
%%%%%%%%%%%%%%%%%%%%%%%%%%%%%%%%%%%%%%%%%%%%%%%%%%%%%%%%%%%%%%%%%%%
\subsection{Fuerza Capilar (\texorpdfstring{$FC$}{FC})}
\label{subsec:fuerza-capilar}

Todo usuario del transporte metropolitano reconoce el fenómeno
aunque no lo nombre: el nodo de Barranca del Muerto, el complejo de
Chapultepec, los grandes paraderos sobre Tlalpan. Son puntos donde
múltiples líneas y sistemas convergen, donde la red se aprieta antes
de abrirse de nuevo hacia las colonias periféricas. La \textbf{Fuerza
Capilar} ($FC$) es el indicador que formaliza esa intuición urbana:
cuantifica el nivel de convergencia de flujos en cada nodo o
agrupación de nodos del grafo, identificando los puntos donde la red
concentra y redistribuye la mayor cantidad de conexiones. El
indicador se evalúa en dos perspectivas complementarias: la
\textbf{Fuerza Capilar de Macro-Hub} ($FC_H$), que agrupa estaciones
físicamente colindantes en un radio de 100~m para representar los
grandes nodos de transbordo, y la \textbf{Fuerza Capilar individual}
($FC$), que mide la convergencia topológica estricta en cada nodo sin
tolerancia espacial.

La expresión formal de la perspectiva de Macro-Hub es:

\begin{equation}
    FC_{H} = \sum_{j \in H}
    \bigl(k_{\text{in}}(j) + k_{\text{out}}(j)\bigr)
    - \text{Aristas Internas}
    \label{eq:fc-geografica}
\end{equation}

\noindent donde:
\begin{itemize}
    \item $H$: conjunto de nodos del Macro-Hub (estaciones agrupadas
          en un radio de tolerancia de 100~m).
    \item $k_{\text{in}}(j)$: número de aristas entrantes del nodo $j$
          dentro del grafo de servicios.
    \item $k_{\text{out}}(j)$: número de aristas salientes del nodo $j$.
    \item Aristas Internas: conexiones entre nodos del propio Hub,
          descontadas para evitar el doble conteo de flujos internos
          al complejo.
\end{itemize}

% ---------------------------------------------------------------
\subsubsection{Distribución por Nodo y Macro-Hub}
\label{subsubsec:fc-distribucion-nodos}
% ---------------------------------------------------------------

La Figura~\ref{fig:grafica-fc-nodos} sintetiza la distribución de la
Fuerza Capilar en los principales nodos de la red, ordenados por su
nivel de concentración de conexiones.

\begin{figure}[H]
    \centering
    \includegraphics[width=0.92\textwidth]{%
        Figures/Mapas/Tableau_Resources/RedActual/VFTModel/UniGrafica/%
        Grafica_VFTModel_2026_Fuerza_capilar}
    \caption[Distribución de Fuerza Capilar en los principales nodos
    de la red]{%
    Distribución de la Fuerza Capilar ($FC$ y $FC_H$) en los
    principales nodos y Macro-Hubs de transferencia de la \zmvm{}.
    Los valores más altos corresponden a los complejos donde mayor
    número de rutas y sistemas se superponen. Corrida: 2026-04-18.}
    \label{fig:grafica-fc-nodos}
    \fuentefigura{Elaboración propia con \textbf{Tableau};
    datos: \texttt{VFTModel} \citep{cabrera2026VFTModel}.}
\end{figure}

% ---------------------------------------------------------------
\subsubsection{Distribución Territorial}
\label{subsubsec:fc-distribucion-territorial}
% ---------------------------------------------------------------

Leída en conjunto con el análisis de cobertura de la sección
anterior, la distribución de la Fuerza Capilar revela la otra cara
del déficit territorial: las demarcaciones con menor cobertura no
carecen simplemente de estaciones, sino de nodos redistributores
capaces de articular los ramales locales con la red troncal. Las
zonas periféricas que registran los menores valores de $C$ son
también las que presentan menor densidad de nodos de alta $FC$: en
ellas, la red no solo es escasa, sino que los pocos nodos existentes
cargan con una proporción desproporcionada de las conexiones locales,
sin el respaldo de un Macro-Hub que diversifique la oferta de
transferencia.

La Figura~\ref{fig:mapa-fuerza-capilar} espacializa esa distribución
sobre la cuenca metropolitana, permitiendo identificar los corredores
y zonas de mayor densidad capilar a escala territorial.

\begin{figure}[H]
    \centering
    \includegraphics[width=0.90\textwidth]{%
        Figures/Mapas/Qgis_Resources/RedActual/Mapa_Fuerza_Capilar}
    \caption[Distribución espacial de la Fuerza Capilar — red
    multimodal \zmvm{}]{%
    Distribución espacial de la Fuerza Capilar ($FC_H$) en la red
    de transporte masivo de la \zmvm{}. Los nodos de mayor
    concentración de conexiones se ubican en el corredor sur de la
    cuenca, sobre los ejes Insurgentes--Tlalpan y el entorno del
    Anillo Periférico. Corrida: 2026-04-18.}
    \label{fig:mapa-fuerza-capilar}
    \fuentefigura{Elaboración propia con \textbf{QGIS};
    datos: \texttt{VFTModel} \citep{cabrera2026VFTModel},
    GTFS \citep{adip_gtfs_2024} e INEGI 2023.}
\end{figure}

% ---------------------------------------------------------------
\subsubsection{Vista Integral de la Red}
\label{subsubsec:fc-vista-integral}
% ---------------------------------------------------------------

La Figura~\ref{fig:dashboard-fuerza-capilar} consolida, mediante el
\termino{dashboard} analítico del modelo, la visión integral del
indicador sobre el conjunto de la red.

\begin{figure}[H]
    \centering
    \includegraphics[width=0.95\textwidth]{%
        Figures/Mapas/Tableau_Resources/RedActual/VFTModel/Dashboards/%
        Dashboard_VFTModel_2026_Fuerza_Capilar}
    \caption[Dashboard de Fuerza Capilar — red multimodal \zmvm{}]{%
    Vista integral de la Fuerza Capilar ($FC$ y $FC_H$) sobre la red
    de transporte masivo de la \zmvm{}, generada con el
    \termino{dashboard} analítico del \texttt{VFTModel}.
    Corrida: 2026-04-18.}
    \label{fig:dashboard-fuerza-capilar}
    \fuentefigura{Elaboración propia con \textbf{Tableau};
    datos: \texttt{VFTModel} \citep{cabrera2026VFTModel}.}
\end{figure}

Los nodos de alta Fuerza Capilar trazan, sobre la ciudad, una
jerarquía implícita que no emergió de una planificación orientada al
transporte colectivo: en una metrópoli cuya infraestructura vial se
estructuró históricamente en función del automóvil privado, la red de
transporte masivo se fue articulando sobre los márgenes e intersticios
de esa trama, y su capacidad redistributiva terminó concentrada en un
puñado de complejos del corredor sur de los que dependen decenas de
rutas periféricas. Esa concentración no es neutral; reproduce, en
términos de conectividad, la misma asimetría territorial documentada
en la cobertura. La sección siguiente evalúa su consecuencia directa
sobre el usuario: el Factor de Desviación mide cuánto se alarga el
viaje cuando esa jerarquía obliga a rodeos.
